\documentclass{article}

\usepackage{hyperref}

\newcommand{\meetingdate}{15 July 2026}
\usepackage[
    a4paper,
    top=3.5cm,
    bottom=2.3cm,
    left=2.3cm,
    right=2.3cm,
    headheight=42pt
]{geometry}
\usepackage[utf8]{inputenc}
\usepackage[T1]{fontenc}
\usepackage{fancyhdr}
\usepackage{lastpage}
\usepackage{enumitem}
\usepackage{amssymb}
\usepackage{etoolbox}
\usepackage{xcolor}
\usepackage{tikz}
\usepackage{tcolorbox}

\newcommand{\project}{Interactive VR Museum}
\newcommand{\meetingtype}{Weekly Meeting}
\providecommand{\meetingdate}{\today}

\newcommand{\authorname}{Jannick Seper}
\newcommand{\tutors}{Rahel Arnold, Raphael Waltenspül}

\lhead{\project\\\meetingtype\\Date: \meetingdate}
\rhead{\authorname\\\tutors\\Page \thepage\ of \pageref{LastPage}}
\pagestyle{fancy}

\newcommand{\done}{
    \section*{1. What I have done}
}

\newcommand{\support}{
    \section*{2. Where I need support}
}

\newcommand{\plan}{
    \section*{3. What I am planning to do for the next meeting}
}

\newtcolorbox{checkpointbox}{
    colback=gray!5,
    colframe=gray!50,
    boxrule=0.5pt,
    arc=3mm,
    left=3mm,
    right=3mm,
    top=2mm,
    bottom=2mm,
    width=\textwidth
}

\newcommand{\checkpoint}[2][open]{
    \begin{checkpointbox}
        \ifstrequal{#1}{done}
            {$\boxtimes$}
            {$\square$}
        \hspace{0.3cm} #2
    \end{checkpointbox}
}

\begin{document}

\done
\begin{itemize}
    \item Prepared separate configurations for standalone deployment and PC streaming, added installation and logging scripts, and documented both workflows.
    \item Improved hand tracking by adding a fallback hand mesh.
    \item Limited the number of loaded results to improve performance.
    \item Improved interaction with displayed media by adding colored grab points.
    \item Moved image and 3D model loading and rendering work to background threads so that the user interface remains responsive during searches.
    \item Cleaned up the codebase using editor rules and removed remaining hard-coded configuration paths.
\end{itemize}

\support
\begin{itemize}
    \item Pre-optimization of 3D models: Can you suggest a method for reducing file sizes without losing too much quality, or a way to remove unsupported textures?
    \item Room with original-size 3D assets: Is this useful? Some assets are scaled to as little as 0.0001\% of their original size, so they would be extremely large in such a room. Should we instead scale each asset to a reasonable size? If so, how should that size be determined, by using minimum and maximum dimensions?
\end{itemize}

\plan
\checkpoint{(p3)Test streaming to the Apple Vision Pro}
\checkpoint{(p1)Investigate how to handle original-size 3D assets}
\checkpoint{(p1)Investigate pre-optimization of 3D models}
\checkpoint{(p2)Add sound effects for clicking and walking}
\checkpoint{(p2)Add a progress bar for asset loading}
\checkpoint{(p1)Prepare for the midterm presentation}

\end{document}
